\documentclass[]{standalone}
\usepackage[T1]{fontenc}\usepackage{tikz}
\usepackage{amsmath, amsfonts}
\usetikzlibrary{arrows, shapes.geometric, calc, positioning, decorations.pathmorphing}

\makeatletter
\def\input@path{{./includes/}{./includes/ut/diags/}}
\makeatother

\begin{document}
\tikzstyle{node} = [rectangle, rounded corners, minimum width=2cm, minimum height=0.8cm, text centered, draw=black]
\tikzstyle{arrow} = [stealth-, shorten >= 2pt, shorten <= 2pt]
\tikzstyle{conflict} = [stealth-stealth, decorate, decoration={zigzag,pre=lineto,pre length=2mm,post=lineto,post length=1mm,segment length=5mm,amplitude=1mm}]
\begin{tikzpicture}[node distance=0.8cm]
    \node (goal) [node] {Safely increase capacity};

    \node (m1) [right = of goal] {};
    \node (mr1) [node, above = of m1 |- goal] {Stay decentralized};
    \node (mr2) [node, below = of m1 |- goal] {Stay secure};
    \draw [arrow] (goal) -- (mr1);
    \draw [arrow] (goal) -- (mr2);

    \node (arp1) [node, right = of mr1.east] {Use many chains and small blocks};
    \node (arp2) [node] at (mr2 -| arp1) {Use one chain and big blocks};
    \draw [arrow] (mr1) -- (arp1);
    \draw [arrow] (mr2) -- (arp2);

    \draw [conflict] (arp1) -- node (c) [anchor=west, inner sep=5pt] {\textbf{Conflict}} (arp2);
    %\node (conflict) [node] at ([xshift=1cm]c.west) {\textbf{Conflict}};
\end{tikzpicture}
\end{document}
